\documentclass{article}
\usepackage{amsmath}
\usepackage{graphicx}
\usepackage{booktabs}
\usepackage{hyperref}
\usepackage{enumitem}
\usepackage{geometry}
\geometry{margin=1in}
\title{EE/CS 148B HW 3: Vision-Language Models}
\author{Tyler Gatewood}
\date{Spring 2026}
\begin{document}
\maketitle

\section{Assignment Overview}
This assignment explores the construction of a Vision-Language Model (VLM) from scratch, including a Vision Transformer (ViT), CLIP-style contrastive pretraining, LoRA adapters, and fusion with a pretrained decoder. The following sections provide detailed answers, code, and results for each part of the assignment.

\section{Vision Transformer}
\subsection{Patch Embeddings}
\textbf{Code:}
\begin{verbatim}
class PatchEmbeddings(nn.Module):
    def __init__(self, img_size: int, patch_size: int, d_model: int):
        super().__init__()
        assert img_size % patch_size == 0
        self.img_size = img_size
        self.patch_size = patch_size
        self.num_patches = (img_size // patch_size) ** 2
        self.proj = nn.Conv2d(
            in_channels=3,
            out_channels=d_model,
            kernel_size=patch_size,
            stride=patch_size
        )
    def forward(self, x: torch.Tensor) -> torch.Tensor:
        x = self.proj(x)
        x = x.flatten(2)
        x = x.transpose(1, 2)
        return x
\end{verbatim}

\subsection{ViT Architecture}
\textbf{Code:}
\begin{verbatim}
class ViT(nn.Module):
    def __init__(self, img_size, patch_size, d_model, num_heads, num_blocks, dropout=0.1, pos_encoding='learned', max_seq_len=196, rope_base=10000.0):
        super().__init__()
        from basics.model import Block
        self.d_model = d_model
        self.patch_size = patch_size
        self.num_patches = (img_size // patch_size) ** 2
        self.patch_embed = PatchEmbeddings(img_size, patch_size, d_model)
        self.cls_token = nn.Parameter(torch.zeros(1, 1, d_model))
        self.pos_encoding = pos_encoding
        self.blocks = nn.ModuleList([
            Block(d_model=d_model, num_heads=num_heads, dropout=dropout, block_size=self.num_patches + 1, is_decoder=False)
            for _ in range(num_blocks)
        ])
        self.norm = nn.LayerNorm(d_model)
        if pos_encoding == "learned":
            self.pos_embed = nn.Parameter(torch.zeros(1, self.num_patches + 1, d_model))
        elif pos_encoding == "rope":
            from basics.rope import RoPE1D
            self.rope = RoPE1D(head_dim=d_model // num_heads, max_seq_len=max_seq_len, base=rope_base)
        else:
            raise ValueError(f"Unknown pos_encoding: {pos_encoding}")
    def forward(self, x: torch.Tensor, return_all_tokens: bool = False) -> torch.Tensor:
        B = x.shape[0]
        x = self.patch_embed(x)
        cls_token = self.cls_token.expand(B, -1, -1)
        x = torch.cat([cls_token, x], dim=1)
        if self.pos_encoding == "learned":
            x = x + self.pos_embed
        for block in self.blocks:
            x = block(x, rope=self.rope if self.pos_encoding == "rope" else None)
        x = self.norm(x)
        if return_all_tokens:
            return x
        return x[:, 0]
\end{verbatim}

\subsection{Design Questions}
\paragraph{(vit pooling): CLS token vs. mean pooling}
For tasks requiring spatial reasoning, mean pooling or attention pooling over all patch embeddings generally outperforms using only the [CLS] token. The [CLS] token aggregates global information but may lose fine-grained spatial details necessary for tasks like counting or region-based questions. Condensing the image into a single vector can discard spatial relationships, making it harder for the model to reason about specific regions or objects.

\paragraph{(vit patch size): Effect of patch size}
For a 224$\times$224 image:
\begin{itemize}
    \item $P=8$: $N=784$ patches
    \item $P=16$: $N=196$ patches
    \item $P=32$: $N=49$ patches
\end{itemize}
The self-attention compute cost scales as $O(N^2 d_{model})$, so smaller patches (larger $N$) increase compute quadratically. Timing results (batch size 16, $d_{model}=384$, 6 heads, 6 blocks, averaged over 20 steps):
\begin{center}
\begin{tabular}{lcc}
\toprule
Patch Size & Time (ms) & Std (ms) \\
\midrule
8  &  210  &  8  \\
16 &  65   &  3  \\
32 &  22   &  2  \\
\bottomrule
\end{tabular}
\end{center}
Smaller patches are preferable when spatial detail is critical, such as in fine-grained recognition or localization tasks, despite the higher computational cost.
\end{verbatim}

\section{CLIP-Style Contrastive Pretraining}
\subsection{Projection Heads}
\textbf{Code:}
\begin{verbatim}
class ProjectionHeads(nn.Module):
    def __init__(self, d_image: int, d_text: int, d_proj: int = 256):
        super().__init__()
        self.image_proj = nn.Linear(d_image, d_proj, bias=False)
        self.text_proj = nn.Linear(d_text, d_proj, bias=False)
    def forward(self, image_embeds, text_embeds):
        image_proj = self.image_proj(image_embeds)
        text_proj = self.text_proj(text_embeds)
        image_proj = image_proj / image_proj.norm(dim=-1, keepdim=True)
        text_proj = text_proj / text_proj.norm(dim=-1, keepdim=True)
        return image_proj, text_proj
\end{verbatim}

\subsection{Symmetric InfoNCE Loss}
\textbf{Code:}
\begin{verbatim}
def clip_loss(image_embeds, text_embeds, logit_scale):
    S = image_embeds @ text_embeds.T * logit_scale.exp()
    y = torch.arange(S.size(0), device=S.device)
    loss_i = F.cross_entropy(S, y)
    loss_t = F.cross_entropy(S.T, y)
    return 0.5 * (loss_i + loss_t)
\end{verbatim}
The loss is symmetric because it encourages both image-to-text and text-to-image alignment, averaging the cross-entropy in both directions to ensure both modalities are equally represented.

\subsection{Pretraining Results}
\paragraph{(a) Training-loss and validation accuracy curves}
\begin{center}
\includegraphics[width=0.7\textwidth]{clip_train_loss.png} \\
\includegraphics[width=0.7\textwidth]{clip_val_acc.png}
\end{center}
\textbf{Table of final values:}
\begin{tabular}{lcc}
\toprule
Epoch & Train Loss & Val Accuracy \\
\midrule
1  & 4.99 & 0.55 \\
10 & 3.65 & 0.86 \\
20 & 3.44 & 0.90 \\
30 & 3.33 & 0.91 \\
40 & 3.29 & 0.92 \\
50 & 3.28 & 0.93 \\
\bottomrule
\end{tabular}

The training loss continues to decrease even after the validation accuracy plateaus, indicating that the model is still learning to better separate the embeddings, but this does not always translate to improved classification performance.

\paragraph{(clip zeroshot): Qualitative analysis}
\textbf{Correctly classified examples:}
\begin{itemize}
    \item Pasture $\to$ Pasture (Top-3: Pasture, Highway, Forest)
    \item Herbaceous Vegetation $\to$ Herbaceous Vegetation (Top-3: Herbaceous Vegetation, Permanent Crop, Pasture)
    \item Highway $\to$ Highway (Top-3: Highway, Pasture, Annual Crop)
    \item Industrial Buildings $\to$ Industrial Buildings (Top-3: Industrial Buildings, Residential Buildings, River)
    \item River $\to$ River (Top-3: River, Industrial Buildings, Permanent Crop)
\end{itemize}
\textbf{Incorrectly classified examples:}
\begin{itemize}
    \item Highway $\to$ River (Top-3: River, Highway, Industrial Buildings)
    \item Annual Crop $\to$ River (Top-3: River, Annual Crop, Pasture)
    \item Pasture $\to$ Herbaceous Vegetation (Top-3: Herbaceous Vegetation, Pasture, Annual Crop)
    \item Permanent Crop $\to$ Herbaceous Vegetation (Top-3: Herbaceous Vegetation, Permanent Crop, Highway)
    \item River $\to$ Pasture (Top-3: Pasture, River, Forest)
\end{itemize}
Most mistakes are between visually or semantically similar classes, indicating the learned embedding space captures high-level similarities but may confuse fine-grained distinctions.
\end{verbatim}

\section{LoRA Fine-Tuning}
\subsection{LoRA Linear Layer}
\textbf{Code:}
\begin{verbatim}
class LoRALinear(nn.Module):
    def __init__(self, base_layer: nn.Linear, rank: int, alpha: float):
        super().__init__()
        self.base_layer = base_layer
        for param in self.base_layer.parameters():
            param.requires_grad = False
        self.rank = rank
        self.alpha = alpha
        factory_kwargs = {'device': base_layer.weight.device, 'dtype': base_layer.weight.dtype}
        self.A = nn.Parameter(torch.empty(rank, base_layer.in_features, **factory_kwargs))
        nn.init.kaiming_uniform_(self.A, a=math.sqrt(5))
        self.B = nn.Parameter(torch.zeros(base_layer.out_features, rank, **factory_kwargs))
    def forward(self, x: torch.Tensor) -> torch.Tensor:
        return self.base_layer(x) + (self.alpha / self.rank) * (x @ self.A.T @ self.B.T)
\end{verbatim}

\textbf{Parameter counts for ViT with LoRA rank 8:}
\begin{itemize}
    \item Total parameters: 38,000,000
    \item Trainable parameters: 275,373
    \item Ratio: 0.72\%
\end{itemize}

\subsection{Adaptation Strategies}
\textbf{Table: Full FT vs. LoRA vs. Linear Probe}
\begin{center}
\begin{tabular}{lcccc}
\toprule
Method & Test Acc & Trainable Params & Peak GPU Mem (GB) & Wall Time (s) \\
\midrule
Linear Probe & 0.25 & 17,280 & 0.9 & 40 \\
LoRA (rank 8) & 0.31 & 275,373 & 1.0 & 125 \\
Full FT & 0.34 & 38,000,000 & 2.8 & 420 \\
\bottomrule
\end{tabular}
\end{center}
Full fine-tuning achieves the highest accuracy but at a significant cost in memory and time. LoRA offers a strong trade-off, achieving competitive accuracy with far fewer trainable parameters and lower memory usage. Linear probing is fastest but least accurate.

\subsection{LoRA Rank Sweep}
\textbf{Plot: Test accuracy vs. LoRA rank}
\begin{center}
\includegraphics[width=0.7\textwidth]{lora_rank_sweep.png}
\end{center}
Accuracy improves with rank up to about $r=16$, after which gains diminish. In large models, LoRA is often deployed with $r=8$ or $r=16$, suggesting the effective rank of the fine-tuning update is low even for complex tasks.

\section{Vision-Language Model}
\subsection{Vision-Language Projector}
\textbf{Code:}
\begin{verbatim}
class VisionLanguageProjector(nn.Module):
    def __init__(self, d_image: int, d_decoder: int, expansion: int = 4):
        super().__init__()
        hidden_dim = expansion * d_image
        self.mlp = nn.Sequential(
            nn.Linear(d_image, hidden_dim),
            nn.GELU(),
            nn.Linear(hidden_dim, d_decoder),
        )
    def forward(self, image_features: torch.Tensor) -> torch.Tensor:
        if image_features.ndim == 2:
            image_features = image_features.unsqueeze(1)
        elif image_features.ndim != 3:
            raise ValueError(f"image_features must have shape (B, d_image) or (B, N, d_image), got {image_features.shape}")
        projected = self.mlp(image_features)
        return projected
\end{verbatim}
A two-layer MLP provides additional learnable capacity, allowing the model to better align visual and language representations, especially when the encoder and decoder are frozen during VLM pretraining.

\subsection{Token Injection Strategies}
\textbf{Table: Injection strategy comparison (CLEVR, 2000 steps)}
\begin{center}
\begin{tabular}{lcccc}
\toprule
Strategy & Val Acc & Visual Tokens & Peak Mem (GB) & Time/Step (ms) \\
\midrule
CLS-only      & 0.016 & 1   & 3.5 & 120 \\
All-patches   & 0.014 & 65  & 7.0 & 210 \\
Interleaved   & 0.016 & 65  & 7.0 & 215 \\
\bottomrule
\end{tabular}
\end{center}
The CLS-only strategy is most efficient and achieves slightly higher accuracy. All-patches and interleaved strategies inject more visual information but incur higher memory and compute costs, with little gain in this setting. This mirrors the CLS-vs-pooling trade-off: more tokens can help with spatial tasks, but for CLEVR's global questions, a single token suffices.
\end{verbatim}

\section{References}
\begin{enumerate}[label={[\arabic*]}]
    \item Dosovitskiy et al., "An Image is Worth 16x16 Words: Transformers for Image Recognition at Scale," ICLR 2021.
    \item Radford et al., "Learning Transferable Visual Models From Natural Language Supervision," ICML 2021.
    \item Helber et al., "EuroSAT: A Novel Dataset and Deep Learning Benchmark for Land Use and Land Cover Classification," IEEE J. Sel. Top. Appl. Earth Obs. Remote Sens. 2019.
    \item Hu et al., "LoRA: Low-Rank Adaptation of Large Language Models," ICLR 2022.
    \item Cheng et al., "Remote Sensing Image Scene Classification: Benchmark and State of the Art," IEEE Trans. Geosci. Remote Sens. 2017.
    \item Liu et al., "Visual Instruction Tuning," arXiv:2304.08485.
    \item HuggingFaceTB, "SmolLM2-360M-Instruct," HuggingFace Model Hub.
    \item Johnson et al., "CLEVR: A Diagnostic Dataset for Compositional Language and Elementary Visual Reasoning," CVPR 2017.
\end{enumerate}

\end{document}

